% -*- coding: utf-8 -*-
% !TEX program = xelatex
\documentclass[plain,noanswer,medmath]{../../template/jnuexam}
\usepackage{../../template/kysx}

\begin{document}


\examtitle{1999}{4}


\loadproblems{1}{1999P1}%载入试卷一
\loadproblems{3}{1999P3}%载入试卷三


\makepart{填空题}{本题共5小题，每小题3分，满分15分}


\begin{problem}
设函数$f(x) = a^x$（$a > 0$, $a \ne 1$），
则$\lim_{n \to \infty} \frac{1}{n^2}\ln[f(1)f(2) \cdots f(n)] =$ \fillin{}．
\end{problem}


\begin{problem}
设$f(x) = \e^x yz^2$，其中$z = z(x,y)$是由$x + y + z + xyz = 0$确定的隐函数，
则$f'_x(0,1, - 1) =$ \fillin{}．
\end{problem}


\useproblem{3}{1}{3}%【同试卷三第一(3)题】


\begin{problem}
已知$AB - B = A$，其中$B = \begin{pmatrix}
  1 & -2 & 0 \\
  2 &  1 & 0 \\
  0 &  0 & 2
\end{pmatrix}$，则$A =$ \fillin{}．
\end{problem}


\begin{problem}
设$X$服从参数为$\lambda$的泊松分布，且已知$E[(X - 1)(X - 2)] = 1$，则$\lambda =$ \fillin{}．
\end{problem}


\makepart{选择题}{本题共5小题，每小题3分，满分15分}


\useproblem{1}{2}{1}%【同试卷一第二(1)题】


\useproblem{3}{2}{2}%【同试卷三第二(2)题】


\useproblem{3}{2}{3}%【同试卷三第二(3)题】


\begin{problem}
设随机变量$X$和$Y$的方差存在且不等于$0$，则$D(X + Y) = DX + DY$是$X$和$Y$\pickout{}
\begin{abcd}
\item 不相关的充分条件，但不是必要条件．
\item 独立的必要条件，但不是充分条件．
\item 不相关的充要条件．
\item 独立的充要条件．
\end{abcd}
\end{problem}


\begin{problem}
设$X$服从指数分布，则$Y = \min \{X,2\}$的分布函数\pickout{}
\begin{abcd}
\item 是连续函数．
\item 至少有两个间断点．
\item 是阶梯函数．
\item 恰有一个间断点．
\end{abcd}
\end{problem}


\makepart{}{本题满分6分}%【同试卷三第三题】

\useproblem{3}{3}{0}


\makepart{}{本题满分7分}%【同试卷三第四题】

\useproblem{3}{4}{0}


\makepart{}{本题满分6分}%【同试卷三第五题】

\useproblem{3}{5}{0}


\makepart{}{本题满分6分}

\begin{problem}
设$F(x)$为$f(x)$的原函数，且当$x \ge 0$时，$f(x)F(x) = \frac{x\e^x}{2(1 + x)^2}$，
已知$F(0) = 1$，$F(x) > 0$，试求$f(x)$．
\end{problem}


\makepart{}{本题满分6分}

\begin{problem}
已知$f(x)$连续，$\int_0^x t f(x - t)\dt = 1 - \cos x$，求$\int_0^{\frac{\pi}{2}} f (x)\dx$的值．
\end{problem}


\makepart{}{本题满分6分}

\begin{problem}
证明：当$0 < x < \pi$时，有$\sin \frac{x}{2} > \frac{x}{\pi}$．
\end{problem}


\makepart{}{本题满分7分}

\begin{problem}
设矩阵$A = \begin{pmatrix}
   3 &  2 & -2 \\
  -k & -1 & k \\
   4 &  2 & -3
\end{pmatrix}$，问：当$k$为何值时，存在可逆矩阵$P$，使得$P^{-1} AP$为对角矩阵？并求出$P$和相应的对角矩阵．
\end{problem}


\makepart{}{本题满分9分}

\begin{problem}
已知线性方程组$\begin{cases}
  x_1 + x_2 + x_3 = 0, \\
  ax_1 + bx_2 + cx_3 = 0, \\
  a^2x_1 + b^2x_2 + c^2x_3 = 0.
\end{cases}$
\begin{enumerate}
\item 当$a$, $b$, $c$满足何种关系时，方程组仅有零解？
\item 当$a$, $b$, $c$满足何种关系时，方程组有无穷多组解，并用基础解系表示全部解．
\end{enumerate}
\end{problem}


\makepart{}{本题满分9分}

\begin{problem}
设二维随机变量$(X,Y)$在矩形$G = \{(x,y)|0 \le x \le 2,0 \le y \le 1\}$上服从均匀分布，
试求边长为$X$和$Y$的矩形面积$S$的概率密度$f(s)$．
\end{problem}


\makepart{}{本题满分8分}

\begin{problem}
已知随机变量$X_1$和$X_2$的概率分布为
$$\begin{array}{|c|cc5c|}
\hline
  X_1 &          -1 &           0 & 1 \\
\hline
    P & \frac{1}{4} & \frac{1}{2} & \frac{1}{4} \\
\hline
\end{array},\qquad \begin{array}{|c|c5c|}
\hline
  X_2 &           0 & 1 \\
\hline
    P & \frac{1}{2} & \frac{1}{2} \\
\hline
\end{array},$$
且$P\{X_1X_2 = 0\} = 1$．
\begin{enumerate}
\item 求$X_1$和$X_2$的联合分布；
\item 问$X_1$和$X_2$是否独立？为什么？
\end{enumerate}
\end{problem}


\end{document}
